\chapter{Related Work}

\section{Overview: Prompt Engineering vs.\ Prompt Optimization}
Prompting is the dominant interface for applying LLMs to downstream tasks without parameter updates. Traditional prompt engineering is largely manual: practitioners iteratively adjust instruction wording, formatting constraints, and examples. While effective, manual tuning is slow and often unstable across models and datasets.

Prompt optimization automates this process by searching over prompts (or prompt components) using an objective such as validation accuracy. The core challenge is cost: each candidate must be evaluated by running an LLM over many examples, so the loop is dominated by LLM calls, tokens, latency, and monetary expense.

\section{Manual and Heuristic Prompt Engineering}
Manual prompt design remains widely used because it is flexible and requires minimal infrastructure. Common heuristics include clear instruction structure, output schemas, few-shot demonstrations, and step-wise decomposition. Its main limitation is scalability: tuning prompts across many tasks/datasets or under strict inference budgets becomes difficult.

\section{Automated Prompt Optimization Methods}
\subsection{Instruction Generation and Selection (APE)}
Automatic Prompt Engineer (APE) uses an LLM to generate candidate instructions and selects high-performing prompts via downstream evaluation~\citep{zhou2022ape}. The approach is simple and effective, but can be expensive when many candidates must be tested.

\subsection{LLM-as-Optimizer (OPRO)}
OPRO treats an LLM as an optimizer that iteratively proposes new candidates based on previously tried prompts and their scores~\citep{yang2023opro}. It demonstrates strong improvements, but repeated evaluation rounds can be costly.

\subsection{Compiler-Based Optimization (DSPy, MIPROv2)}
DSPy provides a programming model for LM pipelines and a compiler that optimizes prompts/demonstrations to maximize a metric~\citep{khattab2023dspy}. MIPROv2 extends this idea with joint optimization of instructions and demonstrations using sample-efficient search strategies~\citep{khattab2024miprov2}. These frameworks are powerful for multi-stage systems, but still rely on repeated evaluations whose cost scales with dataset size and pipeline complexity.

\subsection{Textual ``Gradient'' Optimization (TextGrad)}
TextGrad optimizes textual variables using LLM-generated feedback in an automatic-differentiation-like manner~\citep{yuksekgonul2024textgrad}. While general, it may require multiple LLM calls per update step, making cost a key concern in larger systems.

\section{Evolutionary, Bandit, and RL-Style Search}
Several methods treat prompts as a population and apply mutation/selection (evolutionary search) or allocate limited evaluation budget via exploration--exploitation (bandit-style search). A recurring practical theme is \emph{multi-fidelity evaluation}: cheaply filter many candidates, then spend expensive validation budget only on a small subset.

\section{Cost-Aware Evaluation and Systems Optimizations}
Because LLM inference dominates cost, many systems rely on caching/memoization, batching, and early stopping/adaptive evaluation. Another common technique is surrogate modeling, where a cheap predictor ranks candidates before expensive validation. These optimizations are complementary to the choice of optimizer and often determine real-world feasibility.

\section{Benchmarks}
We use common reasoning/prompting benchmarks: BIG-Bench Hard (BBH)~\citep{suzgun2022bbh}, SciEval~\citep{sun2023scieval}, and MATH (with the MATH-500 subset)~\citep{hendrycks2021math}. These datasets provide diverse task formats and difficulty levels, making them suitable for evaluating both quality and cost trade-offs.

\section{Summary and Connection to This Work}
Prior work shows prompt optimization can yield large gains but is often constrained by evaluation cost. This motivates cost-aware strategies such as surrogate-guided filtering and multi-fidelity evaluation. Our approach follows this direction by using a surrogate model to rank mutations and validating only a small Top-$K$ subset, while enforcing monotonic tracking of the best validation score.